\section{Formaciones Basadas en Estructura Virtual}

Zhen y colaboradores~\cite{zhen2022formation} conciben la formación como un cuerpo rígido virtual al que los vehículos se adhieren, en lugar de perseguir a un líder físico:

\begin{center}
\begin{tikzpicture}[
  font=\fontfamily{ptm}\selectfont\scriptsize,
  >={Stealth[length=2mm]},
  block/.style={draw, minimum width=1.3cm, minimum height=1.3cm, align=center, inner sep=2pt},
  taken/.style={draw, line width=1pt, minimum width=1.3cm, minimum height=1.3cm, align=center, inner sep=2pt},
]
  \node[block] (P)  {Referencia\\Formación};
  \node[taken, right=1.1cm of P] (EV) {Estructura\\Virtual};
  \node[block, right=1.7cm of EV] (CS) {Control de\\Seguimiento};

  \node[align=center, above=0.7cm of EV] (CP) {Campo Potencial};

  \draw[->] (P)  -- node[above]{$\mathbf{p}_r$} node[below]{$[\mathbf{p}_r,\mathbf{u}_r]$} (EV);
  \draw[->, line width=1pt] (EV) -- node[above, pos=0.65]{$\mathbf{p}_{r_i}$} node[below, pos=0.65]{$[\mathbf{p}_{r_i},\mathbf{u}_{r_i}]$} (CS);
  \draw[->] (CP) -- node[right]{$F_i$} (EV);

  \draw[dashed] ($(P.north west)+(-0.2,0.2)$) rectangle ($(EV.south east)+(0.2,-0.2)$);

  \coordinate (uEnd) at ($(CS.east)+(0.8,0)$);
  \draw[->] (CS.east) -- node[above]{$\tau$} node[below]{$[\tau]$} (uEnd);
\end{tikzpicture}
\end{center}
\begin{center}
\footnotesize Formación multi-AUV de Zhen et al.: este proyecto retoma solo la estructura virtual
\end{center}
